\chapter{Notations et lecture des formules}
\label{ch:notations}

\begin{lecture}
Chaque formule met en jeu des objets d'une nature précise et une opération qui
les relie. Les tableaux de ce chapitre donnent pour chaque symbole sa nature, sa
lecture orale et son rôle dans le modèle.
\end{lecture}

Ce chapitre fixe le sens des symboles et la manière de les nommer à l'oral. La
nature d'une notation gouverne les opérations qu'elle admet : un scalaire, un
vecteur, une matrice, une fonction et un ensemble se manipulent selon des règles
propres. La typographie du document la signale — vecteurs en gras, ensembles en
lettres calligraphiques, paramètres grecs en gras lorsqu'ils forment un vecteur.

Le document place l'indice d'observation en bas, selon la convention statistique :
$\x_i$ est le vecteur de l'élève $i$, $x_{ij}$ sa caractéristique $j$ et $c_i$ sa
classe. Certains cours de machine learning écrivent $\x^{(i)}$ pour séparer le
numéro d'exemple d'une puissance ; l'écriture $x_{ij}$ retenue ici s'accorde
directement avec l'usage matriciel et se dicte plus simplement. Les crochets en
exposant, comme $\thv^{[t]}$, portent le numéro d'itération.

\begin{figure}[htbp]
\centering
\resizebox{\textwidth}{!}{%
\begin{tikzpicture}[
  n/.style={draw=accent,fill=accentlight,rounded corners,align=center,
    minimum height=9mm,text width=27mm},
  a/.style={-{Latex},thick,draw=ink},node distance=6mm]
  \node[n] (raw) {$r_{ij}$\\valeur brute};
  \node[n,right=of raw] (x) {$\x_i$\\vecteur préparé};
  \node[n,right=of x] (z) {$z_k$\\score réel};
  \node[n,right=of z] (q) {$q_k$\\sortie binaire};
  \node[n,right=of q] (c) {$\widehat c$\\classe prédite};
  \draw[a] (raw)--node[above,scale=.78]{standardiser}(x);
  \draw[a] (x)--node[above,scale=.78]{$\thv_k^\top\xt$}(z);
  \draw[a] (z)--node[above,scale=.78]{$\sigma$}(q);
  \draw[a] (q)--node[above,scale=.78]{$\Argmax_k$}(c);
\end{tikzpicture}}
\caption{Nature des objets successifs et transformation qui mène de l'un à
l'autre : une valeur brute, un vecteur préparé, un score réel, une sortie
logistique, une classe.}
\label{fig:notation-chain}
\end{figure}

\section{Indices, effectifs et ensembles}

\begin{table}[htbp]
\centering
\begin{tabularx}{\textwidth}{>{$}l<{$} >{\raggedright\arraybackslash}p{3.1cm} X}
\toprule
\text{symbole} & \text{lecture orale} & \text{signification exacte} \\
\midrule
i & « indice i » lorsqu'il est attaché à un symbole & indice d'une observation : ici, un élève. Il varie de $1$ à $m$. \\
j & « indice j » lorsqu'il est attaché à un symbole & indice d'une caractéristique : ici, un cours ou une note. Il varie de $1$ à $n$. \\
k & « indice k » lorsqu'il est attaché à un symbole & indice d'une classe : ici, une maison. Il varie de $1$ à $K$. \\
t & « itération t » & indice d'une itération de l'algorithme d'optimisation. \\
m & « m » & nombre total d'observations dans l'ensemble considéré. \\
n & « n » & nombre de caractéristiques utilisées par le modèle, avant ajout du biais. \\
K & « K » & nombre de classes ; dans ce problème, $K=4$. \\
\mathcal C & « l'ensemble des classes C » & ensemble des classes possibles, et non une classe particulière. On précise « C calligraphique » seulement si l'on doit dicter la graphie. \\
\R & « R » ou « l'ensemble des réels » & ensemble dans lequel prennent leurs valeurs les scores, paramètres et caractéristiques numériques. \\
\mathcal I_{\mathrm{tr}} & « I indice apprentissage » & ensemble des indices appartenant à l'échantillon d'apprentissage. \\
\mathcal I_{\mathrm{te}} & « I indice test » & ensemble des indices appartenant à l'échantillon de test. \\
\bottomrule
\end{tabularx}
\caption{Indices, tailles et ensembles. Un même indice conserve le même rôle dans
tout le document.}
\end{table}

L'écriture $i\in\mathcal I_{\mathrm{tr}}$ se lit « $i$ appartient à l'ensemble
des indices d'apprentissage ». L'écriture $k\in\{1,\ldots,K\}$ se lit « $k$ varie
de un à $K$ ».

\section{Une observation et ses caractéristiques}

\begin{table}[htbp]
\centering
\begin{tabularx}{\textwidth}{>{$}l<{$} >{\raggedright\arraybackslash}p{3.4cm} X}
\toprule
\text{symbole} & \text{lecture orale} & \text{signification exacte} \\
\midrule
r_{ij} & « r indice i j » & valeur brute de la caractéristique $j$ pour l'élève $i$, avant standardisation. Le premier indice repère la ligne, le second la variable. \\
\mu_j & « mu indice j » & moyenne de la caractéristique $j$, estimée uniquement sur l'apprentissage. \\
s_j & « s indice j » & écart-type de la caractéristique $j$, estimé uniquement sur l'apprentissage. \\
x_{ij} & « x indice i j » & caractéristique $j$ standardisée de l'observation $i$ ; c'est un scalaire. \\
\x_i & « x indice i » & vecteur colonne des $n$ caractéristiques standardisées de l'observation $i$ ; $\x_i\in\R^n$. Le caractère gras est généralement implicite à l'oral. \\
\xt_i & « x tilde indice i » ou « x i augmenté » & vecteur $\x_i$ précédé d'une coordonnée égale à $1$ pour intégrer le biais ; $\xt_i\in\R^{n+1}$. \\
\X & « la matrice X » ou « la matrice de design » & matrice qui empile les vecteurs augmentés en lignes ; $\X\in\R^{m\times(n+1)}$. \\
X_i & « la variable aléatoire X indice i » & vecteur aléatoire représentant les caractéristiques de l'observation $i$ avant leur réalisation. Il se distingue de la valeur observée $\x_i$ et de la matrice $\X$. \\
\bottomrule
\end{tabularx}
\caption{Les minuscules désignent une observation ou une composante ; la
majuscule grasse $\X$ réunit toutes les observations sous forme matricielle.}
\end{table}

La formule
\[
  x_{ij}=\frac{r_{ij}-\mu_j}{s_j}
\]
se lit : « x indice i j égale r indice i j moins mu indice j, sur s indice j ».
Elle exprime la valeur de la note en nombre d'écarts-types au-dessus ou au-dessous
de la moyenne d'apprentissage.

\section{Classes observées, cibles binaires et prédictions}

\begin{table}[htbp]
\centering
\begin{tabularx}{\textwidth}{>{$}l<{$} >{\raggedright\arraybackslash}p{3.5cm} X}
\toprule
\text{symbole} & \text{lecture orale} & \text{signification exacte} \\
\midrule
C_i & « la variable aléatoire C indice i » & maison de l'élève $i$ avant observation ; $C_i$ prend ses valeurs dans $\mathcal C$. \\
c_i & « c indice i » & maison effectivement observée pour l'élève $i$ ; $c_i\in\mathcal C$. La minuscule désigne ici une réalisation de $C_i$. \\
y_{ik} & « y indice i k » & cible binaire du classifieur $k$ pour l'observation $i$ : $1$ si $c_i=k$, sinon $0$. \\
\y_k & « le vecteur y indice k » & vecteur colonne réunissant les $m$ cibles du problème « classe $k$ contre le reste » ; $\y_k\in\{0,1\}^m$. \\
\widehat y_i & « y chapeau indice i » & classe binaire prédite pour l'observation $i$, égale à $0$ ou $1$. Le chapeau marque une valeur estimée. \\
\widehat c(\x) & « c chapeau évalué en x » ou « la classe prédite en x » & règle de classification multiclasses évaluée pour les caractéristiques $\x$. \\
\ind[A] & « l'indicatrice de l'événement A » & vaut $1$ lorsque la proposition $A$ est vraie et $0$ sinon. \\
\bottomrule
\end{tabularx}
\caption{La classe observée $c_i$ prend ses valeurs dans $\mathcal C$ ; la cible
binaire $y_{ik}$, construite pour le sous-problème OvR numéro $k$, vaut $0$ ou
$1$.}
\end{table}

L'écriture $y_{ik}=\ind[c_i=k]$ se lit : « y indice i k est l'indicatrice de
l'événement c indice i égale k » ; en contexte, on dira plus naturellement :
« la cible de l'observation i pour la classe k vaut un si sa classe est k ».

\section{Paramètres, score et probabilité binaire}

\begin{table}[htbp]
\centering
\begin{tabularx}{\textwidth}{>{$}l<{$} >{\raggedright\arraybackslash}p{3.7cm} X}
\toprule
\text{symbole} & \text{lecture orale} & \text{signification exacte} \\
\midrule
\theta_j & « thêta indice j » & coefficient scalaire associé à la caractéristique $j$. Son signe indique le sens de son effet sur les log-cotes, toutes choses égales par ailleurs. \\
\theta_0 & « thêta zéro » & biais ou constante du modèle ; il multiplie la coordonnée ajoutée $x_0=1$. \\
\thv & « thêta » ou « le vecteur des paramètres thêta » & vecteur colonne de tous les paramètres $(\theta_0,\ldots,\theta_n)^\top\in\R^{n+1}$. On ne dit « vecteur » que lorsque sa nature doit être soulignée. \\
\thv_k & « thêta indice k » & vecteur des paramètres du classifieur binaire « classe $k$ contre le reste ». \\
z_{\thv}(\x) & « le score z, paramétré par thêta, évalué en x » ; souvent « le score en x » & score linéaire réel $\thv^\top\xt$, à valeurs dans $\R$ tout entier. \\
\sigma & « sigma » & fonction logistique $z\mapsto(1+e^{-z})^{-1}$, qui transforme un score réel en nombre strictement compris entre $0$ et $1$. \\
p_{\thv}(\x) & « p paramétré par thêta, évalué en x » ; souvent « la probabilité prédite en x » & probabilité conditionnelle modélisée de la classe binaire positive : $\Pp(Y=1\mid X=\x)$. \\
q_k(\x) & « q indice k évalué en x » & sortie du classifieur binaire $k$, égale à $\sigma(\thv_k^\top\xt)$. La somme des $K$ valeurs OvR prend une valeur libre. \\
\Theta & « thêta majuscule » ou « la matrice des paramètres » & matrice regroupant les $K$ vecteurs de paramètres, un par classifieur OvR. \\
\bottomrule
\end{tabularx}
\caption{Le modèle produit successivement trois objets : le score $z$ dans $\R$,
la probabilité binaire $p$ dans $]0,1[$, la décision $\widehat y$ dans
$\{0,1\}$.}
\end{table}

Le produit $\thv^\top\xt$ se lit « thêta transposé x tilde », ou « le produit
scalaire de thêta et x tilde ». Le symbole
$\top$ indique la transposition : il transforme ici le vecteur colonne $\thv$ en
vecteur ligne afin que le produit scalaire soit défini. Développé, ce produit vaut
\[
  \thv^\top\xt=\theta_0+\theta_1x_1+\cdots+\theta_nx_n.
\]

\section{Vraisemblance, perte et optimisation}

\begin{table}[htbp]
\centering
\begin{tabularx}{\textwidth}{>{$}l<{$} >{\raggedright\arraybackslash}p{3.7cm} X}
\toprule
\text{symbole} & \text{lecture orale} & \text{signification exacte} \\
\midrule
L(\thv) & « la vraisemblance évaluée en thêta » & vraisemblance : probabilité des étiquettes observées, considérée comme fonction des paramètres. \\
\log L(\thv) & « la log-vraisemblance en thêta » & logarithme de la vraisemblance ; il remplace les produits par des sommes sans changer le maximiseur. \\
\ell_i(\thv) & « ell indice i évaluée en thêta » ou « la perte de l'observation i » & opposé de la contribution de l'observation $i$ à la log-vraisemblance. \\
J(\thv) & « J évaluée en thêta » ou « le coût en thêta » & fonction objectif : moyenne des pertes sur l'échantillon d'apprentissage. \\
\nabla J(\thv) & « le gradient de J en thêta » & vecteur des $n+1$ dérivées partielles de $J$ ; il pointe dans la direction locale de plus forte augmentation. « Nabla J » sert à dicter la formule. \\
\nabla^2J(\thv) & « la Hessienne de J en thêta » & matrice $(n+1)\times(n+1)$ des dérivées secondes ; elle décrit la courbure locale de la fonction objectif. « Nabla carré J » nomme la notation, « Hessienne » nomme l'objet. \\
W & « la matrice W » & matrice diagonale dont le terme $i$ vaut $p_i(1-p_i)$ ; elle intervient dans la Hessienne. \\
\alpha & « alpha » & taux d'apprentissage, c'est-à-dire longueur multiplicative du pas de descente de gradient. \\
\lambda & « lambda » & intensité de la pénalisation quadratique ; $\lambda=0$ signifie absence de régularisation. \\
\eps & « epsilon » & tolérance numérique utilisée dans le critère d'arrêt. \\
\bottomrule
\end{tabularx}
\caption{Objets employés pour estimer les paramètres. La vraisemblance est
maximisée ; son opposé logarithmique $J$ est minimisé.}
\end{table}

La formule
\[
  \nabla J(\thv)=\frac1m\X^\top\bigl(\sigma(\X\thv)-\y\bigr)
\]
se lit : « le gradient de J en thêta vaut un sur m, X transposé, multiplié par
le vecteur sigma de X thêta moins y ». Le vecteur entre parenthèses contient les
écarts entre probabilités prédites et étiquettes observées.

La mise à jour
\[
  \thv^{[t+1]}=\thv^{[t]}-\alpha\nabla J(\thv^{[t]})
\]
se lit : « à l'itération t plus un, thêta vaut sa valeur à l'itération t moins
alpha fois le gradient évalué en cette valeur ». Les crochets $[t]$ repèrent le
numéro d'itération.

\section{Probabilités et opérateurs}

\begin{table}[htbp]
\centering
\begin{tabularx}{\textwidth}{>{$}l<{$} >{\raggedright\arraybackslash}p{3.9cm} X}
\toprule
\text{symbole} & \text{lecture orale} & \text{signification exacte} \\
\midrule
\Pp(A) & « probabilité de l'événement A » & probabilité de l'événement $A$. \\
\Pp(A\mid B) & « probabilité de A conditionnellement à B » ; couramment « A sachant B » & probabilité conditionnelle de $A$ lorsque l'information $B$ est connue. \\
Y\mid X=\x & « la loi de Y conditionnellement à X égal x » & loi conditionnelle de la réponse pour une observation dont les caractéristiques valent $\x$. \\
\sum_{i=1}^m & « somme sur i, de un à m » & addition de l'expression située à droite pour toutes les observations. \\
\prod_{i=1}^m & « produit sur i, de un à m » & multiplication de l'expression située à droite pour toutes les observations. \\
\Argmax_k a_k & « l'argument du maximum sur k de a indice k » & indice $k$ pour lequel $a_k$ atteint son maximum ; le résultat est cet indice, quand $\max_k a_k$ désigne la valeur atteinte. Dans l'usage, « arg max » est également courant. \\
\|v\|_2 & « norme euclidienne de v » ou « norme deux de v » & longueur euclidienne du vecteur $v$. \\
\operatorname{diag}(a_i) & « la matrice diagonale de termes a indice i » & matrice dont la diagonale contient les $a_i$ et dont les autres termes sont nuls. \\
\partial J/\partial\theta_j & « dérivée partielle de J par rapport à thêta indice j » & variation locale de $J$ lorsque seul $\theta_j$ varie, les autres paramètres étant maintenus fixes. \\
\bottomrule
\end{tabularx}
\caption{Opérateurs qui apparaissent dans les dérivations et les règles de
décision.}
\end{table}

Enfin,
\[
  \widehat c(\x)=\Argmax_{k\in\{1,\ldots,K\}}\thv_k^\top\xt
\]
se lit : « la classe prédite en x est l'argument qui maximise, sur les classes k,
le score thêta indice k transposé x tilde ». Cette formule renvoie le nom de la
classe gagnante après conversion de l'indice vers le libellé correspondant.
